% OpenStax Calculus Volume 2 -- Section 7.3 Exercises: Polar Coordinates
% Exercises reproduced from OpenStax, licensed under CC BY 4.0.
% Source: https://openstax.org/books/calculus-volume-3/pages/1-3-polar-coordinates
\documentclass{exercises}
\title{OpenStax Calculus Volume 2}
\subtitle{Section 7.3 Exercises: Polar Coordinates}
\sourceurl{https://openstax.org/books/calculus-volume-3/pages/1-3-polar-coordinates}
\author{}
\date{}

\begin{document}
\maketitle


In the following exercises, plot the point whose polar coordinates are given by first constructing the angle $\theta$ and then marking off the distance $r$ along the ray.

\begin{enumerate}[start=1]
\item $(3,\frac{\pi}{6})$
\item $(-2,\frac{5\pi}{3})$
\item $(0,\frac{7\pi}{6})$
\item $(-4,\frac{3\pi}{4})$
\item $(1,\frac{\pi}{4})$
\item $(2,\frac{5\pi}{6})$
\item $(1,\frac{\pi}{2})$
\end{enumerate}

For the following exercises, consider the polar graph below. Give two sets of polar coordinates for each point.


\textit{[Figure: The polar coordinate plane is divided into 12 sectors. Point A is drawn on the first circle on the first spoke above the $\theta=0$ line in the first quadrant. Point B is drawn in the fourth quadrant on the third circle and the second spoke below the $\theta=0$ line. Point C is drawn on the $\theta=\pi$ line on the third circle. Point D is drawn on the fourth circle on the first spoke below the $\theta=\pi$ line.]}

\begin{enumerate}[resume]
\item Coordinates of point A.
\item Coordinates of point B.
\item Coordinates of point C.
\item Coordinates of point D.
\end{enumerate}

For the following exercises, the rectangular coordinates of a point are given. Find two sets of polar coordinates for the point with angles in $(0,2\pi]$. Round to three decimal places.

\begin{enumerate}[resume]
\item $(2,\,2)$
\item $(3,-4)$
\item $(8,\,15)$
\item $(-6,\,8)$
\item $(4,\,3)$
\item $(3,-\sqrt{3})$
\end{enumerate}

For the following exercises, find rectangular coordinates for the given point in polar coordinates.

\begin{enumerate}[resume]
\item $(2,\frac{5\pi}{4})$
\item $(-2,\frac{\pi}{6})$
\item $(5,\frac{\pi}{3})$
\item $(1,\frac{7\pi}{6})$
\item $(-3,\frac{3\pi}{4})$
\item $(0,\frac{\pi}{2})$
\item $(-4.5,6.5)$
\end{enumerate}

For the following exercises, determine whether the graph of each polar equation is symmetric with respect to the $x$-axis, the $y$-axis, or the origin.

\begin{enumerate}[resume]
\item $r=3\sin(2\theta)$
\item $r^{2}=9\cos \theta$
\item $r=\cos(\frac{\theta}{5})$
\item $r=2\sec \theta$
\item $r=1+\cos \theta$
\end{enumerate}

For the following exercises, describe the graph of each polar equation. Confirm each description by converting into a rectangular equation.

\begin{enumerate}[resume]
\item $r=3$
\item $\theta=\frac{\pi}{4}$
\item $r=\sec \theta$
\item $r=\csc \theta$
\end{enumerate}

For the following exercises, convert the rectangular equation to polar form and sketch its graph.

\begin{enumerate}[resume]
\item $x^{2}+y^{2}=16$
\item $x^{2}-y^{2}=16$
\item $x=8$
\end{enumerate}

For the following exercises, convert the rectangular equation to polar form and sketch its graph.

\begin{enumerate}[resume]
\item $3x-y=2$
\item $y^{2}=4x$
\end{enumerate}

For the following exercises, convert the polar equation to rectangular form and sketch its graph.

\begin{enumerate}[resume]
\item $r=4\sin \theta$
\item $r=6\cos \theta$
\item $r=\theta$
\item $r=\cot \theta \csc \theta$
\end{enumerate}

For the following exercises, sketch a graph of the polar equation and identify any symmetry.

\begin{enumerate}[resume]
\item $r=1+\sin \theta$
\item $r=3-2\cos \theta$
\item $r=2-2\sin \theta$
\item $r=5-4\sin \theta$
\item $r=3\cos(2\theta)$
\item $r=3\sin(2\theta)$
\item $r=2\cos(3\theta)$
\item $r=3\cos(\frac{\theta}{2})$
\item $r^{2}=4\cos(2\theta)$
\item $r^{2}=4\sin \theta$
\item $r=2\theta$
\item \techrequired The graph of $r=2\cos(2\theta)\sec(\theta)$ is called a strophoid. Use a graphing utility to sketch the graph and, from the graph, determine the asymptote.
\item \techrequired Use a graphing utility and sketch the graph of $r=\frac{6}{2\sin \theta-3\cos \theta}$.
\item \techrequired Use a graphing utility to graph $r=\frac{1}{1-\cos \theta}$.
\item \techrequired Use technology to graph $r=e^{\sin(\theta)}-2\cos(4\theta)$.
\item \techrequired Use technology to plot $r=\sin(\frac{3\theta}{7})$ (use the interval $0\le \theta \le 14\pi$).
\item Without using technology, sketch the polar curve $\theta=\frac{2\pi}{3}$.
\item \techrequired Use a graphing utility to plot $r=\theta \sin \theta$ for $-\pi \le \theta \le \pi$.
\item \techrequired Use technology to plot $r=e^{-0.1\theta}$ for $-10\le \theta \le 10$.
\item \techrequired There is a curve known as the ``Black Hole.'' Use technology to plot $r=e^{-0.01\theta}$ for $-100\le \theta \le 100$.
\item \techrequired Use the results of the preceding two problems to explore the graphs of $r=e^{-0.001\theta}$ and $r=e^{-0.0001\theta}$ for $|\theta|>100$.
\end{enumerate}

\end{document}
